\chapter{2026:Jacob Tsimerman}\label{ch2026jacobtsimerman:jacob-tsimerman-fields-medal-2026}

\begin{center}\rule{0.5\linewidth}{0.5pt}\end{center}

\textbf{Prepared:} July 2026 \textbf{Subject:} Jacob Tsimerman (b. April 26, 1988, Kazan, Russia) \textbf{Position:} Professor of Mathematics, University of Toronto \textbf{Doctoral Advisor:} Peter Sarnak (Princeton University) \textbf{Fields Medal Citation:} \emph{``For his work on the Andr\'{e}--Oort conjecture and for his contributions to Hodge theory, including the development of o-minimal GAGA.''}

\begin{center}\rule{0.5\linewidth}{0.5pt}\end{center}

\begin{prereqs} Algebraic geometry (varieties, moduli spaces), number theory (Galois theory, arithmetic of elliptic curves), basic complex analysis, model theory. \end{prereqs}

\begin{objectives} State the Andr\'{e}--Oort conjecture and explain its significance; describe the strategy of combining model theory with arithmetic geometry; understand the role of o-minimal structures in algebraic geometry. \end{objectives}

\section{Problem History}

Shimura varieties are higher-dimensional geometric spaces that parameterize abelian varieties with additional structure. They arise naturally in the Langlands program, which seeks deep connections between number theory, algebraic geometry, and representation theory. For example, the modular curve $Y(N)$ parametrizes elliptic curves with level-$N$ structure, and higher-dimensional Shimura varieties parametrize abelian varieties with additional structures such as polarization, level structure, and endomorphisms.

Special points on Shimura varieties correspond to abelian varieties with particularly rich arithmetic properties, such as complex multiplication (CM). A CM abelian variety has endomorphism ring larger than $\mathbb{Z}$, giving it extraordinary symmetry. The Andr\'{e}--Oort conjecture, posed by Yves Andr\'{e} in 1989 and Frans Oort in 1995, states: \emph{If a closed irreducible algebraic subset of a Shimura variety contains a dense set of special points, then it is a Shimura subvariety.}

The conjecture is a vast generalization of results in transcendental number theory. For curves in modular varieties, it reduces to the Schanuel conjecture and results on linear forms in logarithms. For higher-dimensional Shimura varieties, it describes the algebraic structure of sets determined purely by their arithmetic content.

The conjecture had been proved in special cases: for Hilbert modular surfaces (Tsimerman, 2012), for products of modular curves (Pila--Tsimerman, 2014), and in cases of ``large Galois orbits'' established by Pila and others. But the full conjecture for arbitrary Shimura varieties---particularly for the moduli space $\mathcal{A}_g$ of principally polarized abelian varieties of dimension $g$---had been a central open problem for over two decades.

The \textbf{Pila--Wilkie theorem} (2006) is a landmark result in model theory that provides bounds on the number of rational points on definable sets that lie outside a finite union of algebraic curves. The theorem states that for a definable set $X \subset \mathbb{R}^n$ and $\epsilon > 0$, the number of rational points of height at most $H$ in $X \setminus Z$, where $Z$ is a finite union of algebraic curves, is at most $C_\epsilon H^\epsilon$. The theorem has become a standard tool in arithmetic geometry through Tsimerman's applications, and it provides the bridge between model theory and Diophantine geometry that is essential for the proof of the Andr\'{e}--Oort conjecture.

The Pila--Wilkie theorem is proved using tools from model theory, including cell decomposition and quantifier elimination. The key idea is to decompose the definable set $X$ into cells where the set is ``tame''---meaning it has controlled geometry and can be approximated by algebraic sets. The number of rational points in each cell is then bounded using the theory of heights and the properties of the cell decomposition. The exceptional set $Z$ consists of algebraic curves that are ``too well-approximated'' by rational points, and the theorem shows that these curves can be removed without affecting the bound.

The \textbf{Galois orbit bounds} established by Tsimerman in his PhD thesis are essential for the proof. For a special point $x$ on a Shimura variety, the Galois orbit $\mathrm{Gal}(\overline{\mathbb{Q}}/\mathbb{Q}) \cdot x$ has size at least $c \cdot D(x)^{1/2}$, where $D(x)$ is the discriminant of the number field generated by $x$. This bound shows that special points are ``spread out'' in a quantifiable way, which is essential for the counting argument. The proof uses the theory of complex multiplication and the properties of Shimura varieties to obtain the lower bound on the Galois orbit size.

The Galois orbit of a special point $x$ is the set of all points obtained by applying the symmetries of the number field in which $x$ is defined. For a CM abelian variety $A$, the Galois orbit consists of all CM abelian varieties that are conjugate to $A$ under the action of the Galois group. The size of the Galois orbit is related to the degree of the number field generated by the endomorphisms of $A$, which is controlled by the discriminant of the number field.

The lower bound on the Galois orbit size is essential for the counting argument: it shows that special points are ``spread out'' across the Shimura variety, which means that a subvariety containing a dense set of special points must have large dimension. This is the key step in the proof of the Andr\'{e}--Oort conjecture: if a subvariety contains a dense set of special points, then the Galois orbits must be large, which forces the subvariety to be a Shimura subvariety.

The proof of the Andr\'{e}--Oort conjecture proceeds by contradiction. Suppose there exists a Shimura variety $X$ and a closed irreducible subvariety $Z$ that contains a dense set of special points but is not itself a Shimura subvariety. The Galois orbit bound shows that the special points in $Z$ are spread out across $X$, which means that $Z$ must have large dimension. The Pila--Wilkie theorem then provides bounds on the number of rational points on definable sets that lie outside a finite union of algebraic curves, which contradicts the assumed density of special points in $Z$.

The concept of ``complexity'' for a special point is defined using the discriminant of the associated number field. A special point on a Shimura variety corresponds to an abelian variety with complex multiplication, and its complexity is measured by the size of the discriminant of the CM field. The Galois orbit bound shows that as the complexity increases, the size of the Galois orbit increases at least polynomially. This polynomial growth is sufficient for the counting argument in the Pila--Wilkie method.

Tsimerman's earlier work established lower bounds on the sizes of Galois orbits of special points, showing that special points must be ``spread out'' in a quantifiable way. The Galois orbit of a special point is the set of all points obtained by applying the symmetries of the number field in which the point is defined. These bounds were combined with the Pila--Wilkie theorem---a landmark result in model theory providing bounds on rational points on definable sets outside algebraic curves---to prove important special cases.

A separate but related problem concerned Hodge theory and Griffiths' conjecture: whether the intermediate Jacobian of a projective algebraic variety can be realized as an algebraic group. This question connects analytic and algebraic geometry in a fundamental way.

\section{How Was It Solved}

In 2021, Tsimerman, Peter Pila, and Ananth Shankar posted a proof of the full Andr\'{e}--Oort conjecture. The proof proceeds by contradiction: suppose there exists a Shimura variety $X$ and a closed irreducible subvariety $Z$ that contains a dense set of special points but is not itself a Shimura subvariety.

The strategy has three key ingredients. First, effective bounds on Galois orbits of special points show that special points cannot cluster too tightly. If $Z$ contains a dense set of special points, the Galois orbits must be large, meaning the points are spread out across $X$ in a controlled way. Tsimerman's PhD thesis established these bounds, showing that the size of the Galois orbit of a special point is bounded below by a power of its ``complexity'' (measured by the discriminant of the associated number field).

Second, the Pila--Wilkie theorem provides bounds on the number of rational points on definable sets that lie outside a finite union of algebraic curves. By showing that the set of special points (or related algebraic points) can be described within an o-minimal structure, one can apply this theorem to obtain counting estimates that contradict the assumed density of special points in $Z$.

Third, Tsimerman developed new results on the structure of definable sets in o-minimal structures arising from Shimura varieties. The key innovation was to show that certain sets have hidden algebraic structure, allowing the Pila--Wilkie method to be applied in settings where previous approaches failed.

The proof of Griffiths' conjecture (2018, with Benjamin Bakker and Yohan Brunebarbe) introduced the framework of \textbf{o-minimal GAGA}: a generalization of Serre's classical GAGA principle (relating algebraic and analytic geometry) to o-minimal structures. An \emph{o-minimal structure} is a collection of definable sets and functions satisfying tameness conditions ensuring sets lack pathological behavior---no infinite discrete sets, no wild oscillations.

In o-minimal GAGA, one works with spaces defined by functions in an o-minimal structure rather than by algebraic or analytic functions. The framework shows that if such a space satisfies certain compactness and regularity conditions, then it must be algebraic. This was used to prove that the intermediate Jacobian of a projective variety is an algebraic group, resolving Griffiths' conjecture.

The intermediate Jacobian $J(V)$ of a projective algebraic variety $V$ of dimension $n$ is a complex torus defined using the Hodge structure on the cohomology of $V$. Griffiths conjectured that $J(V)$ is always an algebraic group, meaning it has a group structure compatible with its analytic structure. The proof uses o-minimal GAGA to show that the intermediate Jacobian can be realized as a quotient of an o-minimal group, and the o-minimal GAGA principle then implies that this quotient is algebraic.

An \textbf{o-minimal structure} on $\mathbb{R}$ is a collection $\mathcal{S}$ of subsets of $\mathbb{R}^n$ (for all $n$) satisfying: (1) $\mathcal{S}$ contains all algebraic sets; (2) $\mathcal{S}$ is closed under finite Boolean operations and projections; (3) every set in $\mathcal{S}$ is a finite union of cells. The key property is that every definable set in an o-minimal structure has a well-behaved cell decomposition, which ensures that the set lacks pathological behavior such as infinite discrete subsets or wild oscillations.

The \textbf{GAGA principle} (Serre, 1956) states that for a projective algebraic variety $X$ over $\mathbb{C}$, the categories of algebraic coherent sheaves and analytic coherent sheaves on $X$ are equivalent. This principle shows that algebraic and analytic geometry are essentially the same for projective varieties. O-minimal GAGA extends this principle to o-minimal structures, showing that for a definable space satisfying certain conditions, the categories of algebraic and definable coherent sheaves are equivalent.

The o-minimal GAGA principle has implications for the study of periods of algebraic varieties. The period matrix of an abelian variety is a $g \times g$ complex matrix that encodes the Hodge structure, and the algebraic relations between periods are described by the Andr\'{e}--Oort conjecture. O-minimal GAGA provides a framework for understanding these relations by showing that the period matrix can be realized as a definable function in an o-minimal structure.

\section{Impact of This Discovery}

The Andr\'{e}--Oort conjecture provides a classification of algebraic subvarieties of Shimura varieties that are ``determined'' by their arithmetic content. This classification has implications for the study of special values of $L$-functions, the distribution of prime numbers, and the Langlands program.

The proof demonstrated the power of combining techniques from model theory, arithmetic geometry, and analytic number theory. The Pila--Wilkie theorem, originally a purely model-theoretic result, has become a standard tool in arithmetic geometry through Tsimerman's applications.

The proof of the Andr\'{e}--Oort conjecture has implications for the study of special values of $L$-functions. The $L$-function of an abelian variety encodes information about its arithmetic properties, and the Andr\'{e}--Oort conjecture describes the algebraic relations between special points on Shimura varieties that correspond to abelian varieties with complex multiplication. Understanding these relations has implications for the study of special values of $L$-functions and the theory of automorphic forms.

O-minimal GAGA has proven to be a versatile framework, with applications to periods of algebraic varieties, transcendence properties of special values, and the structure of Hodge loci---the subsets of moduli spaces where a Hodge class exists.

Tsimerman's work connects to the deep arithmetic of abelian varieties. The special points correspond to abelian varieties with complex multiplication, and the Andr\'{e}--Oort conjecture describes when the ``density'' of such varieties forces the ambient space to have arithmetic structure.

The techniques developed have influenced the study of unlikely intersections---the study of when algebraic subvarieties contain ``too many'' points of special type---connecting to the Mordell--Lang conjecture, the Manin--Mumford conjecture, and Zilber's Pink conjectures.

The \textbf{Shimura varieties} are higher-dimensional algebraic varieties that parametrize abelian varieties with additional structure. For example, the modular curve $Y(N)$ parametrizes elliptic curves with level-$N$ structure, and the Siegel modular variety $\mathcal{A}_g$ parametrizes principally polarized abelian varieties of dimension $g$. The special points on a Shimura variety correspond to abelian varieties with complex multiplication, which are abelian varieties with endomorphism ring larger than $\mathbb{Z}$. The Andr\'{e}--Oort conjecture describes the algebraic structure of sets determined by their arithmetic content, and its proof has implications for the study of special values of $L$-functions and the Langlands program.

The \textbf{Langlands program} is a vast web of conjectures connecting number theory, algebraic geometry, and representation theory. The Andr\'{e}--Oort conjecture is related to the Langlands program through the theory of automorphic forms: special points on Shimura varieties correspond to automorphic forms with complex multiplication, and the Andr\'{e}--Oort conjecture describes the algebraic relations between these automorphic forms. The proof of the Andr\'{e}--Oort conjecture provides evidence for the Langlands program and suggests new directions for research.

\section{Modern Views and Improvements}

Tsimerman's o-minimal GAGA framework continues to find new applications. Recent work has used it to study the transcendence properties of periods, the structure of variations of Hodge structure, and new cases of the Andr\'{e}--Oort conjecture for more general Shimura varieties and Shimura stacks.

The connection between model theory and arithmetic geometry has been further developed. New applications of the Pila--Wilkie method to Diophantine geometry, including results on integral points on algebraic varieties and the distribution of rational points, build on Tsimerman's techniques.

The Pila--Wilkie method has been extended to other settings, including the study of integral points on algebraic varieties (where the goal is to understand when algebraic varieties have ``too many'' integral points) and the distribution of rational points on curves and surfaces (where the goal is to understand the density of rational points). These applications have led to new results on the Mordell conjecture (Faltings' theorem) and on the distribution of rational points on elliptic curves.

Tsimerman has also been an outspoken voice on the impact of artificial intelligence on mathematics, arguing that the field must adapt to changing technological realities while preserving the depth and creativity that define mathematical research.

Future directions include extending the Andr\'{e}--Oort conjecture to more general Shimura varieties, developing further applications of o-minimal structures to algebraic geometry, and exploring connections to the theory of periods and motivic Galois groups.

The connection between the Andr\'{e}--Oort conjecture and the theory of \textbf{unlikely intersections} is an active area. The Zilber--Pink conjecture predicts that algebraic subvarieties of Shimura varieties that contain ``too many'' special points must be contained in proper Shimura subvarieties or in special subvarieties defined by additional endomorphisms. The Andr\'{e}--Oort conjecture is a special case of this broader framework, and the techniques developed by Tsimerman have been applied to other cases of the Zilber--Pink conjecture.

The theory of \textbf{periods} of algebraic varieties connects Tsimerman's work to transcendence theory. The period matrix of an abelian variety is a $g \times g$ complex matrix that encodes the Hodge structure, and the algebraic relations between periods are described by the Andr\'{e}--Oort conjecture. Understanding these relations has implications for the study of special values of L-functions and the theory of automorphic forms.

The development of o-minimal structures over non-archimedean fields is another direction. Recent work has extended the Pila--Wilkie theorem to p-adic fields, which has potential applications to the Andr\'{e}--Oort conjecture over local fields. The challenge is to develop the theory of o-minimal structures in the p-adic setting, where the topology is fundamentally different from the real topology.

The Langlands program provides a broad framework for understanding the arithmetic of Shimura varieties. The proof of the Andr\'{e}--Oort conjecture has implications for the Langlands program by describing the algebraic relations between special points that correspond to automorphic forms with complex multiplication. Understanding these relations has implications for the construction of Galois representations and the study of their properties.

The study of special values of $L$-functions is closely connected to the Andr\'{e}--Oort conjecture. The $L$-function of an abelian variety encodes information about its arithmetic properties, and the special values of $L$-functions are related to the periods of the abelian variety. The Andr\'{e}--Oort conjecture describes the algebraic relations between special points on Shimura varieties, which correspond to abelian varieties with complex multiplication. These algebraic relations have implications for the study of special values of $L$-functions and for the Birch and Swinnerton-Dyer conjecture.

The o-minimal GAGA framework has been applied to the study of Hodge loci. The Hodge locus of a variation of Hodge structure is the set of points where a particular Hodge class exists, and understanding the algebraic structure of this set is a major problem in algebraic geometry. The o-minimal GAGA principle provides a framework for studying Hodge loci by showing that they have tameness properties that follow from the o-minimal structure. This approach has led to new results on the algebraic structure of Hodge loci and on the transcendence properties of periods.

The connection between Tsimerman's work and the theory of unlikely intersections has been further developed. The Zilber-Pink conjecture, which generalizes the Andr\'{e}-Oort conjecture, predicts that algebraic subvarieties of Shimura varieties that contain "too many" special points must be contained in proper Shimura subvarieties or in special subvarieties defined by additional endomorphisms. The techniques developed by Tsimerman for the Andr\'{e}-Oort conjecture have been applied to other cases of the Zilber-Pink conjecture, including results on the intersection of algebraic subvarieties with Shimura subvarieties in products of modular curves.

The development of the Pila-Wilkie method for p-adic fields is another active direction. The p-adic version of the Pila-Wilkie theorem would have applications to the study of integral points on p-adic varieties and to the proof of p-adic analogues of the Andr\'{e}-Oort conjecture. The challenge is to develop the theory of o-minimal structures in the p-adic setting, where the topology is fundamentally different from the real topology and the tools of model theory need to be adapted to the new setting.

The study of Galois representations associated to abelian varieties is closely connected to the Andr\'{e}-Oort conjecture. The Galois representation on the Tate module of an abelian variety encodes information about its arithmetic properties, and the Andr\'{e}-Oort conjecture describes the algebraic relations between special points that correspond to abelian varieties with complex multiplication. Understanding these relations has implications for the construction of Galois representations and for the study of their properties.

The p-adic Langlands program provides a framework for studying the p-adic properties of Galois representations. The construction of the p-adic local Langlands correspondence, which relates p-adic Galois representations to p-adic automorphic representations, requires a detailed understanding of the cohomology of p-adic local systems on p-adic analytic spaces. The development of perfectoid spaces by Scholze has provided new tools for studying this correspondence, and the connection between Tsimerman's work on the Andr\'{e}-Oort conjecture and the p-adic Langlands program is an active area of research.

The study of periods of abelian varieties is closely connected to the Andr\'{e}-Oort conjecture. The period matrix of an abelian variety is a complex matrix that encodes the Hodge structure, and the algebraic relations between periods are described by the Andr\'{e}-Oort conjecture. The development of new techniques for studying periods, including the theory of p-adic periods and the theory of motivic periods, is an active area that builds on Tsimerman's foundational work. The connection between periods and the Langlands program, through the theory of automorphic forms and L-functions, is a deep and active area of research.

The study of o-minimal structures over non-archimedean fields has implications for the p-adic version of the Andr\'{e}-Oort conjecture. The p-adic Andr\'{e}-Oort conjecture predicts that algebraic subvarieties of p-adic Shimura varieties that contain too many special points must be contained in proper Shimura subvarieties. The development of o-minimal structures in the p-adic setting, where the topology is fundamentally different from the real topology, is an active area of research that builds on Tsimerman's foundational work on the real version of the conjecture.

\begin{center}\rule{0.5\linewidth}{0.5pt}\end{center}

\section*{Summary}
\addcontentsline{toc}{section}{Summary}

Tsimerman proved the full Andr\'{e}--Oort conjecture using Galois orbit bounds and the Pila--Wilkie method, and developed o-minimal GAGA, bridging model theory and algebraic geometry in profound ways.
